\chapter{Meet the Rust OTA Server}

\section{Read main.rs as a morning routine}

Open \codepath{ota-server/src/main.rs}. The file may look formal at first, yet its job is close to opening a workshop in the morning. Turn on the lights. Unlock the cabinets. Put the tools on the bench. Open the front door.

The Rust version of that routine goes in this order:

\begin{enumerate}
  \item Start logging so the terminal can tell us what happens.
  \item Read settings from \codepath{.env}.
  \item Open the SQLite database.
  \item Prepare the firmware builder.
  \item Prepare the coordinator that watches each build.
  \item Put those shared pieces into \codepath{AppState}.
  \item Connect the HTTP routes.
  \item Listen on the configured address and port.
\end{enumerate}

Here is the middle of that routine with some details folded away:

\begin{lstlisting}[language=Rust]
let config = Arc::new(Config::load()?);
let storage = Storage::connect(&config.database_url()).await?;
let builder = Arc::new(LocalFirmwareBuilder::new(/* settings */));
let coordinator = Arc::new(BuildCoordinator::new(
    storage.clone(),
    builder,
    config.project_dir.clone(),
    config.build_dir.clone(),
));

let state = AppState { config, storage, coordinator };
\end{lstlisting}

The word \codepath{Arc} means several requests can safely hold a shared reference to the same service object. You do not need to master it now. Notice the shape: configuration, storage, and building meet in one state bundle, and route handlers receive that bundle when they need it.

\section{Configuration is the page of house rules}

Move to \codepath{ota-server/src/config.rs}. This file reads each \codepath{OTA_...} setting, supplies local defaults, checks values, and resolves folder paths from the repository root.

It also creates \codepath{firmware-projects}, \codepath{firmware-builds}, and \codepath{data} when they are missing. That is why a fresh copy of the project can start without you building those folders by hand.

If no public URL is supplied, the server tries to detect a likely LAN address. Detection is convenient on a simple home network. An explicit \codepath{OTA_PUBLIC_BASE_URL} is easier to reason about when a PC has Ethernet, Wi-Fi, virtual adapters, or a VPN.

When startup succeeds, the terminal prints two views:

\begin{lstlisting}[numbers=none]
OTA Server running
Local:
http://localhost:7000

LAN / public firmware URL:
http://192.168.1.5:7000
\end{lstlisting}

The local line is for the browser on the PC. The LAN line is the address an ESP32 can reach.

\section{Routes are the labels on the front desk}

Open \codepath{ota-server/src/routes/mod.rs}. This file is short because it mainly matches an address with the function that handles it.

The work is grouped into three nearby files:

\begin{description}
  \item[\codepath{routes/firmware.rs}] Receives project ZIP files and serves finished firmware.
  \item[\codepath{routes/builds.rs}] Starts builds and returns build history or logs.
  \item[\codepath{routes/devices.rs}] Remembers boards, heartbeats, deployments, update checks, and OTA progress.
\end{description}

Find this line:

\begin{lstlisting}[language=Rust,numbers=none]
.route("/api/ota/status", get(status))
\end{lstlisting}

Read it from left to right: when a GET request arrives at that address, call \codepath{status}. Other lines follow the same grammar.

\section{Why the browser is allowed through}

The page comes from port 5000 and calls a server on port 7000. Browsers treat those as different origins. Axum therefore needs a CORS rule that names the dashboard origins it trusts.

The default list contains:

\begin{lstlisting}[numbers=none]
http://localhost:5000
http://127.0.0.1:5000
\end{lstlisting}

If you serve the frontend from another address, add that exact origin to \codepath{OTA_ALLOWED_ORIGINS}. A CORS error in the browser can happen even while the Axum status URL opens normally in its own tab.

\begin{warningbox}
Keep the origin list narrow. Allowing every website to call a local firmware builder turns a browser convenience into a needless open door.
\end{warningbox}

\section{Errors should speak in complete sentences}

Open \codepath{ota-server/src/errors.rs}. Every ordinary API error carries an HTTP status, a short machine-readable code, and a message for a person. It may also carry details and a hint.

\begin{lstlisting}[language=JSON,numbers=none]
{
  "success": false,
  "error": {
    "code": "INVALID_FIRMWARE_PROJECT",
    "message": "Managed application entrypoint was not found",
    "details": "src/main.rs"
  }
}
\end{lstlisting}

The frontend can show the message without guessing what went wrong. A developer can use the short code in tests. The status number gives HTTP clients the broad category:

\begin{description}
  \item[400] The request could not be read.
  \item[404] The named record does not exist.
  \item[409] The request clashes with current state, such as a second active build.
  \item[413] The upload is too large.
  \item[422] The request was readable but failed a project rule.
  \item[500] Something inside the server failed unexpectedly.
\end{description}

\section{SQLite is the server's notebook}

The database file is \codepath{data/ota.db}. You do not need to create its tables. \codepath{storage.rs} does that during startup.

Each table answers a different question:

\begin{longtable}{p{0.19\textwidth}p{0.7\textwidth}}
\toprule
Table & The question it answers\\
\midrule
\endhead
projects & Which ZIP was uploaded, for which target and version?\\
builds & What happened when we tried to compile it?\\
firmware & Which successful binary can a device download?\\
devices & Which boards have introduced themselves, and what did they report?\\
deployments & Which firmware was assigned to which board, and how far did it get?\\
\bottomrule
\end{longtable}

SQLite foreign keys tie those pages together. A firmware file that still belongs to a device or deployment cannot be casually deleted. The API returns a conflict instead of tearing a page out of the notebook and leaving dangling references behind.

\begin{checkpoint}
Open \codepath{main.rs}, \codepath{config.rs}, \codepath{routes/mod.rs}, and \codepath{storage.rs} in four tabs. In each tab, find the one place that matches its chapter role: startup, settings, addresses, and database tables.
\end{checkpoint}

